\documentclass[12pt]{article}
\usepackage{amsmath, amssymb}
\usepackage{geometry}
\geometry{margin=1in}
\begin{document}

\section*{Gaussian Random Variable and Associated Functions}


\section*{1. Moments and Central Moments}

\subsection*{1.1 Definition: Central Moment}

Let $X$ be a continuous random variable with PDF $f_X(x)$ and mean $\mu_X = \mathbb{E}[X]$.
  
The $n^{\text{th}}$ central moment is defined as:

\[
\mathbb{E}\left[(X - \mu_X)^n\right] = \int (x - \mu_X)^n f_X(x)\, dx
\]

\textbf{Dependencies explicitly used:}
\begin{itemize}
\item Expectation is defined via integration over the PDF.
\item The mean $\mu_X$ must exist.
\item The expression depends only on $f_X(x)$ and $\mu_X$.
\end{itemize}

\subsection*{1.2 Evaluation of Central Moments for Specific Orders}

\textbf{Case $n=0$}

\[
\mathbb{E}[(X - \mu_X)^0] = \mathbb{E}[1] = 1
\]

This follows directly from the definition of expectation.

\textbf{Case $n=1$}

\[
\mathbb{E}[X - \mu_X] = \mathbb{E}[X] - \mu_X = \mu_X - \mu_X = 0
\]

\textbf{Reasoning used:}

Linearity of expectation and the definition of the mean.

\textbf{Case $n=2$: Variance}

\[
\sigma_X^2 = \mathbb{E}\left[(X - \mu_X)^2\right]
\]

Expand the square:

\[
(X - \mu_X)^2 = X^2 - 2\mu_X X + \mu_X^2
\]

Taking expectation:

\[
\mathbb{E}[X^2] - 2\mu_X \mathbb{E}[X] + \mu_X^2
\]

Substituting $\mathbb{E}[X] = \mu_X$:

\[
\sigma_X^2 = \mathbb{E}[X^2] - \mu_X^2
\]

\subsection*{1.3 Higher-Order Central Moments}

\textbf{Skewness ($n=3$)}

\[
C_s = \frac{\mathbb{E}\left[(X - \mu_X)^3\right]}{\sigma_X^3}
\]

\textbf{Lecture interpretation:}
\begin{itemize}
\item $C_s > 0$: PDF is right-skewed
\item $C_s < 0$: PDF is left-skewed
\end{itemize}

\textbf{Kurtosis ($n=4$)}

\[
C_k = \mathbb{E}\left[(X - \mu_X)^4\right]
\]

\textbf{Lecture statement:}

A large value implies the random variable has a large peak near the mean.

\section*{2. Properties of Expectation}

\subsection*{2.1 Linearity with Constants}

For constants $a$ and $b$:

\[
\mathbb{E}[aX + b] = a\mathbb{E}[X] + b
\]

\textbf{Dependency:}

Follows directly from linearity of integration.

\subsection*{2.2 Linearity over Sums}

If a function is written as:

\[
g(x) = g_1(x) + g_2(x) + \cdots + g_N(x)
\]

Then:

\[
\mathbb{E}\left[\sum_{k=1}^{N} g_k(X)\right] = \sum_{k=1}^{N} \mathbb{E}[g_k(X)]
\]

\subsection*{2.3 Expectation via PDF}

\[
\mathbb{E}[X] = \int x f_X(x)\, dx
\]

This is the foundational definition used in all derivations.

\section*{3. Gaussian Random Variable}

\subsection*{3.1 Definition}

A random variable $X$ is Gaussian if its PDF is:

\[
f_X(x) = \frac{1}{\sqrt{2\pi\sigma^2}} \exp\left(-\frac{(x-m)^2}{2\sigma^2}\right)
\]

Where:
\begin{itemize}
\item $m = \mathbb{E}[X]$
\item $\sigma^2 = \mathrm{Var}(X)$
\end{itemize}

Notation:

\[
X \sim \mathcal{N}(m,\sigma^2)
\]

\subsection*{3.2 Standard Normal Distribution}

If:

\[
m = 0, \quad \sigma^2 = 1
\]

Then the distribution is the standard normal distribution.

\section*{4. Standard Functions Associated with Gaussian RVs}

\subsection*{4.1 Error Function}

\[
\mathrm{erf}(x) = \frac{2}{\sqrt{\pi}} \int_{0}^{x} e^{-t^2}\, dt
\]

\subsection*{4.2 Complementary Error Function}

\[
\mathrm{erfc}(x) = 1 - \mathrm{erf}(x) = \frac{2}{\sqrt{\pi}} \int_{x}^{\infty} e^{-t^2}\, dt
\]

\subsection*{4.3 $\Phi$-Function (CDF of Standard Normal)}

\[
\Phi(x) = \frac{1}{\sqrt{2\pi}} \int_{-\infty}^{x} e^{-t^2/2}\, dt
\]

\subsection*{4.4 Q-Function (Gaussian Tail Function)}

\[
Q(x) = \frac{1}{\sqrt{2\pi}} \int_{x}^{\infty} e^{-t^2/2}\, dt
\]

Identity stated in lecture:

\[
Q(x) = 1 - \Phi(x)
\]

\section*{5. Relation Between $\Phi$-Function and Q-Function}

\subsection*{5.1 Evaluating the CDF of a Gaussian RV}

Start from the CDF definition:

\[
F_X(x) = \int_{-\infty}^{x} \frac{1}{\sqrt{2\pi\sigma^2}} \exp\left(-\frac{(y-m)^2}{2\sigma^2}\right)\, dy
\]

Substitute:

\[
t = \frac{y-m}{\sigma}
\]

Then:

\[
F_X(x) = \Phi\left(\frac{x-m}{\sigma}\right)
\]

\subsection*{5.2 Evaluating Tail Probabilities}

\[
\Pr(X > x) = \int_{x}^{\infty} \frac{1}{\sqrt{2\pi\sigma^2}} \exp\left(-\frac{(y-m)^2}{2\sigma^2}\right)\, dy
\]

After substitution:

\[
\Pr(X > x) = Q\left(\frac{x-m}{\sigma}\right)
\]

\subsection*{5.3 Key Identities (Explicitly Listed)}

\[
Q(x) = 1 - \Phi(x)
\]

\[
F_X(x) = 1 - Q\left(\frac{x-m}{\sigma}\right)
\]

\textbf{Lecture notes:}
\begin{itemize}
\item $\Phi$-function is used to evaluate the CDF
\item Q-function is more natural for probabilities of the form $\Pr(X > x)$
\end{itemize}

\section*{6. Gaussian Example (Fully from Lecture)}

\textbf{Given PDF}

\[
f_X(x) = \frac{1}{\sqrt{8\pi}} \exp\left(-\frac{(x+3)^2}{8}\right)
\]

\subsection*{Step 1: Identify Parameters}

Compare with the Gaussian PDF:

\[
\frac{1}{\sqrt{2\pi\sigma^2}} \exp\left(-\frac{(x-m)^2}{2\sigma^2}\right)
\]

Thus:

\[
m = -3, \quad \sigma^2 = 4, \quad \sigma = 2
\]

\subsection*{Step 2: Gaussian Relations Used}

\[
F_X(x) = \Phi\left(\frac{x-m}{\sigma}\right)
\]

\[
\Pr(X > x) = Q\left(\frac{x-m}{\sigma}\right)
\]

\subsection*{Step 3: Probabilities to Be Evaluated}

\begin{enumerate}
\item $\Pr(X \le 0)$

Left-tail probability $\rightarrow$ $\Phi$-function

\item $\Pr(X > 4)$

Right-tail probability $\rightarrow$ Q-function

\item $\Pr(|X + 3| < 2)$

Convert:

\[
-2 < X + 3 < 2 \Rightarrow -5 < X < -1
\]

Evaluate as:

\[
F_X(-1) - F_X(-5)
\]

\item $\Pr(|X - 2| > 1)$

Converted into tail probabilities and evaluated using $\Phi$ and Q as appropriate
\end{enumerate}

\textbf{Lecture hint explicitly stated:}

Use Q(x) for right-tail probabilities and $\Phi(x)$ for left-tail probabilities.

\end{document}
